\par
\SourceLeaf{082}
Aus der Elementaralgebra ist der Zusammenhang zwischen den Wurzeln der Gleichung $X=0$ und den einfachen Factoren von $X$ bekannt. Geschieht nemlich
jener Gleichung durch die Substitution $x=p$ Genüge, so ist $x-p$ ein Factor
von $X$, und gibt es $n$ verschiedene Arten, jener Gleichung Genüge zu leisten,
nemlich durch $x=p$, $x=p'$, $x=p''$ u.\,s.\,w., so wird das Product $(x-p)$
$(x-p')(x-p'')\, .\, .\, .\, .$ mit $X$ identisch sein. Unter besondern Umständen kann
aber auch eine Auflösung, wie $x=p$, in $X$ den Factor $(x-p)^2$, oder $(x-p)^3$
oder irgend eine höhere Potenz bedingen, in welchen Fällen man die Wurzel $p$
wie zweimal, dreimal u.\,s.\,w. vorhanden betrachtet.

Verlangt man also nur den Beweis, dass die Function $X$ gewiss \textit{einen} einfachen Factor zulasse, so ist es zureichend, nur das Vorhandensein irgend einer
Wurzel der Gleichung $X=0$ nachzuweisen. Soll aber die vollständige Zerlegbarkeit der Function in einfache Factoren auf Einmal bewiesen werden, so muss
gezeigt werden, dass der Gleichung $X=0$ Genüge geleistet werden kann, entweder durch $n$ ungleiche Werthe von $x$, oder durch eine zwar geringere Anzahl
ungleicher Auflösungen, wovon aber ein Theil die Charactere der mehrfach geltenden gleichen Wurzeln dergestalt an sich trägt, dass die Zusammenzählung aller ungleichen und gleichen die Totalsumme $=n$ hervorbringt.

\FifthArticle{2}
Das ganze Gebiet der complexen Grössen, in welchem die der Gleichung
$X=0$ genügenden Werthe von $x$ gesucht werden sollen, ist ein Unendliches
von zwei Dimensionen, indem, wenn ein solcher Werth $x=t+iu$ gesetzt wird
(wo $i$ immer die imaginäre Einheit $\surd{-1}$ bedeutet), für $t$ und $u$ alle reellen
Werthe von $-\infty$ bis $+\infty$ zulässig sind. Wir haben nun zuvörderst aus diesem unendlichen Gebiete ein abgegrenztes endliches auszuscheiden, ausserhalb
dessen gewiss keine Wurzel der bestimmten Gleichung $X=0$ liegen kann.
Dies kann auf mehr als Eine Art geschehen: unserm Zweck am meisten gemäss
scheint die folgende zu sein.

Anstatt der Form $t+iu$ gebrauche man diese
\[
x=r(\cos\rho+i\sin\rho)
\]
wonach zur Umfassung des ganzen unendlichen Gebiets der complexen Grössen
$r$ durch alle positiven Werthe von $0$ bis $+\infty$, und $\rho$ von $0$ bis $360^0$, oder, was
\SourceLeaf{083}
dasselbe ist, von einem beliebigen Anfangswerthe bis an einen um $360^0$ grössern
Endwerth ausgedehnt werden muss.

Um für $r$ eine Grenze zu, erhalten, über welche hinaus kein Werth mehr
einer Wurzel der Gleichung $X=0$ entsprechen kann, setze ich zuvörderst die
Coëfficienten der einzelnen Glieder von $X$ in eine ähnliche Form, wie $x$, nemlich
\[
\begin{array}{l}
A=a(\cos\alpha+i\sin\alpha)\\
B=b(\cos\beta+i\sin\beta)\\
C=c(\cos\gamma+i\sin\gamma)\quad\mbox{u.\,s.\,w.}
\end{array}
\]
wo also $a,b,c$ bestimmte positive Grössen bedeuten sollen, abgesehen davon,
dass auch eine oder die andere darunter $=0$ sein kann. Ich betrachte sodann
die Gleichung
\[
r^n-\surd2\,.(ar^{n-1}+br^{n-2}+cr^{n-3}+\mbox{u.\,s.\,w.})=0
\]
welche, wie man leicht sieht, eine positive Wurzel hat, und zwar (\textsc{Harriots} Lehrsatz zufolge) nur Eine solche. Es sei $R$ diese Wurzel, wo dann von selbst klar
ist, dass für jeden positiven Werth von $r$, der grösser ist als $R$, der Werth von
$r^n-\surd2\,.(ar^{n-1}+br^{n-2}+cr^{n-3}+\mbox{u.\,s.\,w.})$ positiv sein, und dass dasselbe auch
von der Function
\[
nr^n-\surd2\,.((n-1)ar^{n-1}+(n-2)br^{n-2}+(n-3)cr^{n-3}+\mbox{u.\,s.\,w.})
\]
gelten wird, da dieselbe das $n$fache der ersten Function um
\[
\surd2\,.(ar^{n-1}+2br^{n-2}+3cr^{n-3}+\mbox{u.\,s.\,w.})
\]
also um eine positive Differenz übertrifft.

\FifthArticle{3}
Ich behaupte nun, dass die Grösse $R$ geeignet ist, eine solche Grenze für
die Werthe von $r$, wie im vorhergehenden Artikel gefordert ist, abzugeben. Der
Beweis dieses Satzes ist auf folgende Art zu führen.

Ich setze allgemein $X=T+iU$, wo selbstredend $T$ und $U$ reelle Grössen
bedeuten, und zwar wird
\[
\begin{array}{r@{}l}
T={}&r^n\cos n\rho+ar^{n-1}\cos((n-1)\rho+\alpha)+br^{n-2}\cos((n-2)\rho+\beta)\\
&{}+cr^{n-3}\cos((n-3)\rho+\gamma)+\mbox{u.\,s.\,w.}
\end{array}
\]
\SourceLeaf{084}
\[
\begin{array}{r@{}l}
U={}&r^n\sin n\rho+ar^{n-1}\sin((n-1)\rho+\alpha)+br^{n-2}\sin((n-2)\rho+\beta)\\
&{}+cr^{n-3}\sin((n-3)\rho+\gamma)+\mbox{u.\,s.\,w.}
\end{array}
\]
Man übersieht leicht, dass wenn für $r$ irgend ein positiver Werth grösser als $R$
gewählt wird, $T$ nothwendig dasselbe Zeichen haben wird wie $\cos n\rho$, so oft
dieser Cosinus absolut genommen nicht kleiner ist als $\surd\frac12$. Man braucht nemlich
nur $T$ in folgende Form zu setzen
\[
\begin{array}{r@{}l}
\pm T={}&\surd\frac12\,.r^n-ar^{n-1}-br^{n-2}-cr^{n-3}-\mbox{u.\,s.\,w.}\\
&{}+[\pm\cos n\rho-\surd\frac12]r^n\\
&{}+[1\pm\cos((n-1)\rho+\alpha)]ar^{n-1}\\
&{}+[1\pm\cos((n-2)\rho+\beta)]br^{n-2}\\
&{}+[1\pm\cos((n-3)\rho+\gamma)]cr^{n-3}\\
&{}+\mbox{u.\,s.\,w.}
\end{array}
\]
wo die obern Zeichen für den Fall eines positiven, die untern für den Fall eines
negativen $\cos n\rho$ gelten sollen, und wo der erste Theil des Ausdrucks auf der
rechten Seite positiv ist, in Folge des im vorhergehenden Artikel gegebenen Satzes,
von den folgenden aber wenigstens keiner negativ werden kann. Auf ganz ähnliche Weise erhellet (indem man in obiger Formel nur $U$ anstatt $T$ und durchgehends Sinus anstatt Cosinus schreibt), dass unter gleicher Voraussetzung in Beziehung auf $r$, allemal $U$ dasselbe Zeichen hat wie $\sin n\rho$, so oft dieser Sinus
absolut genommen nicht kleiner ist als $\surd\frac12$. Es hat demnach in allen Fällen wenigstens die eine der beiden Grössen $T,U$ ein voraus bestimmtes positives oder
negatives Zeichen, und es kann folglich für keinen Werth von $\rho$ die Function
$X=0$ werden. W.~Z.~B.~W.

\FifthArticle{4}
Um das Verhalten von $T$ und $U$ in Beziehung auf die Zeichen und deren
Wechsel (bei einem bestimmten, $R$ überschreitenden, Werthe von $r$) noch mehr
ins Licht zu setzen, lasse man $\rho$ alle Werthe zwischen zwei um $360^0$ verschiedenen Grenzen durchlaufen, wozu jedoch nicht $0$ und $360^0$, sondern, indem zur
Abkürzung
\[
\frac{45^0}{n}=\omega
\]
\SourceLeaf{085}
gesetzt wird, $-\omega$ und $(8n-1)\omega$ gewählt werden sollen. Den ganzen Zwischenraum theile ich in $4n$ gleiche Theile, so dass der erste sich von $-\omega$ bis $\omega$,
der zweite von $\omega$ bis $3\omega$, der dritte von $3\omega$ bis $5\omega$ u.\,s.\,w. erstreckt. Zuvörderst hat man auch noch die Werthe der Differentialquotienten $\frac{\mathrm{d}T}{\mathrm{d}\rho}$, $\frac{\mathrm{d}U}{\mathrm{d}\rho}$ in Betracht zu ziehen, wofür man hat
\[
\begin{array}{rl}
\frac{\mathrm{d}T}{\mathrm{d}\rho}
&=-nr^n\sin n\rho-(n-1)ar^{n-1}\sin\bigl((n-1)\rho+\alpha\bigr)\\
&\qquad -(n-2)br^{n-2}\sin\bigl((n-2)\rho+\beta)\bigr)\\
&\qquad -(n-3)cr^{n-3}\sin\bigl((n-3)\rho+\gamma\bigr)-\mathrm{u.\,s.\,w.}\\
\frac{\mathrm{d}U}{\mathrm{d}\rho}
&=nr^n\cos n\rho+(n-1)ar^{n-1}\cos\bigl((n-1)\rho+\alpha\bigr)\\
&\qquad +(n-2)br^{n-2}\cos\bigl((n-2)\rho+\beta)\bigr)\\
&\qquad +(n-3)cr^{n-3}\cos\bigl((n-3)\rho+\gamma\bigr)+\mathrm{u.\,s.\,w.}
\end{array}
\]
Man erkennt daraus leicht, durch ähnliche Schlüsse wie im vorhergehenden Artikel und unter Zuziehung des Satzes am Schlusse von Art. 2, dass $\frac{\mathrm{d}T}{\mathrm{d}\rho}$ immer
das entgegengesetzte Zeichen von $\sin n\rho$ hat, so oft dieser Sinus absolut genommen nicht kleiner ist als $\surd\frac{1}{2}$, dass hingegen $\frac{\mathrm{d}U}{\mathrm{d}\rho}$ immer dasselbe Zeichen wie
$\cos n\rho$ hat, so oft der absolute Werth dieses Cosinus nicht kleiner ist als $\surd\frac{1}{2}$.
Hieraus zieht man folgende Schlüsse.

In dem ersten Intervalle, d.\,i. von $\rho=-\omega$ bis $\rho=+\omega$, ist $T$ stets
positiv, $U$ hingegen für den Anfangswerth negativ, für den Endwerth positiv,
mithin dazwischen gewiss einmal $=0$, und zwar \emph{nur} einmal, weil in dem ganzen Intervalle $\frac{\mathrm{d}U}{\mathrm{d}\rho}$ positiv ist.

In dem zweiten Intervalle ist $U$ stets positiv, $T$ zu Anfang positiv, am
Ende negativ, dazwischen einmal $T=0$ und zwar nur einmal, weil in dem
ganzen Intervalle $\frac{\mathrm{d}T}{\mathrm{d}\rho}$ negativ ist.

In dem dritten Intervalle ist $T$ stets negativ, $U$ einem Zeichenwechsel unterworfen, so dass \emph{einmal} $U=0$ wird.

Im vierten Intervalle ist $U$ stets negativ, $T$ \emph{einmal} $=0$.

In den folgenden Intervallen wiederholen sich in gleicher Ordnung diese
Verhältnisse, so dass das fünfte dem ersten, das sechste dem zweiten u.\,s.\,f.
gleichsteht.

\FifthArticle{5}
Aus der im vorhergehenden Artikel erörterten Folgeordnung der positiven
und negativen Werthe von $T$ und $U$, die bei jedem über $R$ hinausgehenden
\SourceLeaf{086}
Werthe von $r$ Statt findet\HistoricalNote{14}, lässt sich nun folgern, dass innerhalb des Gebiets
der kleinern Werthe von $r$ gewisse Kreuzungen in diesen Anordnungen vorhanden
sein müssen, die das Wesen unsers zu beweisenden Lehrsatzes in sich schliessen.
Ich werde die Beweisführung in einer der Geometrie der Lage entnommenen Einkleidung darstellen, weil jene dadurch die grösste Anschaulichkeit und Einfachheit gewinnt. Im Grunde gehört aber der eigentliche Inhalt der ganzen Argumentation einem höhern von Räumlichem unabhängigen Gebiete der allgemeinen
abstracten Grössenlehre an, dessen Gegenstand die nach der Stetigkeit zusammenhängenden Grössencombinationen sind, einem Gebiete, welches zur Zeit noch
wenig angebaut ist, und in welchem man sich auch nicht bewegen kann ohne
eine von räumlichen Bildern entlehnte Sprache.

\FifthArticle{6}
Das ganze Gebiet der complexen Grössen wird vertreten durch eine unbegrenzte Ebene, in welcher jeder Punkt, dessen Coordinaten in Beziehung auf
zwei einander rechtwinklig schneidende Achsen $t,u$ sind, als der complexen
Grösse $x=t+iu$ entsprechend betrachtet wird: bringt man diese complexe
Grösse in die Form $x=r(\cos\rho+i\sin\rho)$, so bedeuten $r,\rho$ die Polarcoordinaten
des entsprechenden Punkts. Der Inbegriff aller complexen Grössen, für welche
$r$ einerlei bestimmten Werth hat, wird demnach durch einen Kreis repräsentirt,
dessen Halbmesser dieser Werth, und dessen Mittelpunkt der Anfangspunkt der
Coordinaten ist. Denjenigen dieser Kreise, für welchen $r$ um eine nach Belieben gewählte Differenz grösser als $R$ ist, will ich mit $K$ bezeichnen, und mit
$(1),\,(2),\,(3)\,.\,.\,.\,.\,(2n)$ diejenigen Punkte auf demselben, welchen die beziehungsweise zwischen $\omega$ und $3\omega$, zwischen $5\omega$ und $7\omega$, zwischen $9\omega$ und $11\omega$ u.\,s.\,f.
bis zwischen $(8n-3)\omega$ und $(8n-1)\omega$ liegenden Werthe von $\rho$ entsprechen,
für welche nach dem 4. Artikel $T=0$ wird. Man bemerke dabei, dass für die
Punkte $(1),\,(3),\,(5)$ u.\,s.\,w. $U$ positiv, für die Punkte $(2),\,(4),\,(6)$ u.\,s.\,w. hingegen
negativ sein wird.
\SourceLeaf{087}
\FifthArticle{7}
Die Gesammtheit derjenigen Punkte in unserer Ebene, für welche $T$ positiv ist, bildet zusammenhängende Flächentheile, wie schon von selbst erhellet,
wenn man erwägt, dass bei einem stetigen Uebergange von einem Punkte zu einem andern $T$ sich nach der Stetigkeit ändert. Eben so bilden sämmtliche
Punkte, für welche $T$ negativ wird, zusammenhängende Flächentheile. Zwischen den Flächentheilen der ersten Art und denen der zweiten liegen Punkte, in
welchen $T=0$ wird, und nach der Natur der Function $T$ können diese Punkte
nicht auch Flächenstücke, sondern nur Linien bilden, welche einerseits die einen, andererseits die andern Flächentheile begrenzen.

Der ausserhalb $K$ liegende Raum enthält $n$ Flächen der ersten Art, die mit
eben so vielen der zweiten Art abwechseln, und wovon jede, von einem Stück der
Kreislinie $K$ an, zusammenhängend sich ins Unendliche erstreckt. Zugleich aber
ist klar, dass jedes dieser Flächenstücke sich über die Kreislinie hinaus in den
innern Raum fortsetzt, und dass in Beziehung auf die weitere Gestaltung folgende
Fälle Statt finden können.

1) Das betreffende von einem Theile von $K$ anfangende Flächenstück endigt sich isolirt innerhalb der Kreisfläche; seine peripherische Begrenzung besteht
dann nur aus zwei zusammenhängenden Stücken, wovon eines ein Bestandtheil
von $K$ ist, das andere innerhalb des Kreisraumes liegt. In der beigefügten Figur, welche sich auf eine Gleichung fünften Grades bezieht und wo die Zeichen
von $T$ in den verschiedenen Flächentheilen eingeschrieben sind, finden sich drei
der Flächen mit positivem $T$ in diesem Falle; die eine hat die Grenzlinien $10.1$
und $1.11.10$; die zweite diese $4.5$ und $5.12.4$; die dritte $6.7$ und $7.13.6$.
Flächentheile ähnlicher Art mit negativem $T$ finden sich zwei vor.

2) Das Flächenstück durchsetzt einfach die Kreisfläche dergestalt, dass es
mit einem an einer andern Stelle eintretenden Eine zusammenhängende Fläche
bildet. Die ganze peripherische Begrenzungslinie wird dann aus vier Stücken bestehen, von denen zwei der Kreislinie $K$ angehören, und die beiden andern dem
innern Raume. In unserer Figur findet sich dieser Fall bei dem durch $2.3;$
$3.0.8;\ 8.9;\ 9.11.2$ begrenzten Flächenstück.

3) Das Flächenstück spaltet sich im innern Kreisraume einmal oder mehreremale dergestalt, dass es mit noch zweien oder mehrern an andern Stellen eintretenden eine zusammenhängende Fläche bildet, deren ganze peripherische \SourceLeaf{088}Begrenzung dann aus sechs, acht oder mehrern Stücken in gerader Zahl bestehen
wird, die abwechselnd der Kreislinie und dem innern Raume angehören. In
unserer Figur tritt dies ein bei einem Flächentheile, dessen Begrenzung durch
die sechs Stücke $3.4;\ 4.12.5;\ 5.6;\ 6.13.7;\ 7.8;\ 8.0.3$ gebildet wird, in
welchem aber $T$ negativ ist.

\FifthArticle{8}

Bei einer vollständigen Aufzählung aller denkbaren Gestaltungen der in den
innern Kreisraum eintretenden Flächentheile würden den angegebenen Fällen
noch anderweitige Modificationen beigefügt werden müssen. Wenn z. B. ein solcher Flächentheil sich zwar in zwei Aeste spaltet, diese aber im innern Raume
sich wieder vereinigen, so würde dieser Fall, jenachdem nach der Vereinigung
die Fläche im Innern ihren Abschluss findet, oder (ohne neue Theilung) sich bis
zu einer andern Stelle der Kreislinie fortsetzt, dem ersten oder zweiten Falle des
vorhergehenden Artikels zugerechnet werden können, indem die Gestaltung der
Fläche nur durch das Einschliessen einer nicht zu ihr gehörenden Insel modificirt
sein würde. Uebrigens würde es nicht schwer sein, strenge zu beweisen, dass
bei der besondern Beschaffenheit der Function $T$ Modificationen dieser Art \emph{gar
nicht möglich} sind: für unsern Zweck ist dies jedoch unnöthig, indem es nur auf
die Folge der Stücke der \emph{äussern} Begrenzung jedes der in Rede stehenden Flächentheile (d. i. derjenigen, in welchen $T$ positiv ist) ankommt.

Wir haben nemlich schon bemerklich gemacht, dass die Anzahl dieser Stücke
allemal gerade ist (zwei im ersten Falle des vorhergehenden Artikels, vier im
zweiten, sechs oder mehrere im dritten), wovon wechselsweise eines der Kreislinie
$K$, eines dem innern Raume angehört. Ferner ist klar, dass wenn jene äussere
Begrenzungslinie immer in einerlei Sinn durchlaufen wird, wozu hier derjenige
gewählt werden soll, in welchem die Bezifferungen der Punkte von $K$ wachsen
(also, Beispiels halber in unserer Figur so, dass die Fläche immer rechts von der
Begrenzungslinie liegt), der Anfangspunkt und der Endpunkt eines der Kreislinie
angehörenden Stücks beziehungsweise durch eine gerade und die um eine Einheit
grössere ungerade Zahl bezeichnet sein wird, mithin der Anfangspunkt und der
Endpunkt jedes den innern Raum durchlaufenden Stücks allemal beziehungsweise
durch eine ungerade und eine gerade Zahl.

Es steht also fest, dass von den $n$ an einem mit einer ungeraden Zahl \SourceLeaf{089}bezeichneten Punkte von $K$ in den innern Raum eintretenden Linien, in denen
überall $T=0$ ist, eine jede auf eine ganz bestimmte Art\HistoricalNote{15} diesen Raum zusammenhängend durchläuft, bis sie an einer andern mit einer geraden Zahl bezeichneten Stelle wieder austritt. Da nun, wie schon oben (Schluss des 6. Art.) bemerkt
ist, in ihrem Anfangspunkte den Werth von $U$ positiv, am Endpunkte negativ
ist, so muss wegen der Stetigkeit der Werthänderung nothwendig in einem Zwischenpunkte $U=0$ werden. Dieser Punkt repräsentirt dann eine Wurzel der
Gleichung $X=0$; und da die Anzahl solcher Linien $=n$ ist, so ergeben sich
auf diese Weise allemal $n$ Wurzeln jener Gleichung.

\FifthArticle{9}

Wenn die gedachten Linien durch den Kreisraum gehen ohne ein Zusammentreffen mit einander, so ist klar, dass die so erhaltenen $n$ Wurzeln nothwendig ungleich sind. Ein solches freies Durchgehen findet sich in unsrer Figur bei
den Linien von $3$ nach $8$, von $5$ nach $4$ und von $7$ nach $6$, und es gehören
dazu die durch die Punkte $0,\ 12,\ 13$ repräsentirten Wurzeln. Wenn hingegen
zwei solcher Linien, oder mehrere, einen Punkt gemeinschaftlich haben, so ist
zwar darum noch nicht nothwendig, aber doch möglich, dass dieser Punkt zugleich derjenige ist, in welchem $U=0$ wird, in welchem Falle dann zwei oder
mehrere Wurzeln in Eine zusammenfallen, oder, wie es gewöhnlich ausgedrückt
wird, unter sich gleich sein werden. In unsrer Figur treffen die Linien $1.10$
und $9.2$ in dem Punkte $11$ zusammen, und in demselben wird zugleich $U=0$;
die Gleichung hat also ausser den schon aufgeführten drei ungleichen noch zwei
gleiche Wurzeln.

\FifthArticle{10}

Es bleibt nur noch übrig, nachzuweisen, dass wenn der eine Wurzel $=p$
\SourceLeaf{090}
repräsentirende Punkt $P$ in zweien oder mehrern Linien $T=0$ zugleich liegt,
das Quadrat von $x-p$ oder die der Anzahl jener concurrirenden Linien entsprechende höhere Potenz in $X$ als Factor enthalten sein wird. Der Beweis davon
beruhet auf folgenden Sätzen.

Man führe anstatt der unbestimmten Grösse $x$ eine andere $z$ ein, indem
man $x=z+p$ setzt. Es gehe durch diese Substitution $X$ in $Z$ über, wo also
$Z$ eine Function von $z$ von gleicher Ordnung wie $X$ von $x$ sein wird, deren constantes Glied aber fehlt. Indem man dieselbe nach aufsteigenden Potenzen von
$z$ ordnet, sei das niedrigste nicht verschwindende Glied
\[
=Kz^m\quad\mbox{und}\quad Z=Kz^m(1+\zeta)
\]
wo $\zeta$ die Form $Lz+L'zz+L''z^3+\mbox{u.\,s.\,w.}+\frac{1}{K}z^{n-m}$ haben wird; endlich setze
man
\[
z=s(\cos\psi+i\sin\psi)
\]

Der reelle und der imaginäre Bestandtheil von $z$ drücken die Lage jedes
unbestimmten Punkts der Ebene als rechtwinklige Coordinaten, und die Grössen
$s$, $\psi$ die Polarcoordinaten ganz eben so relativ gegen den Punkt $P$ aus, wie die
Bestandtheile von $x$, und die Grössen $r$, $\varphi$ die relative Lage gegen den ursprünglichen Anfangspunkt bezeichnen. Die Verbindung eines bestimmten Werthes
von $s$ mit allen Werthen von $\psi$ in einer Ausdehnung von $360^0$ stellt also die
Punkte einer Kreislinie dar, die ihren Mittelpunkt in $P$ hat und deren Halbmesser $=s$ ist.

Setzt man nun $K=k(\cos\kappa+i\sin\kappa)$, und folglich
\[
Kz^m=ks^m(\cos(m\psi+\kappa)+i\sin(m\psi+\kappa))
\]
so wird für ein unendlich kleines $s$ die Grösse $\zeta$, die wenigstens von derselben
Ordnung ist wie $s$, neben der $1$ vernachlässigt, und mithin gesetzt werden dürfen
\[
T=ks^m\cos(m\psi+\kappa)
\]
woraus erhellet, dass während $\psi$ um $360^0$ wächst, das Zeichen von $T$ in $m$
Stücken der Kreisperipherie positiv, und in eben so vielen mit jenen abwechselnden negativ ist, oder dass $T$ in $2m$ Punkten $=0$ wird, nemlich für
$\psi=\frac{1}{m}(\kappa-90^0)$, $\frac{1}{m}(\kappa+90^0)$, $\frac{1}{m}(\kappa+270^0)$ u.\,s.\,w. Es gehen demnach von $P$ \SourceLeaf{091}zusammen $2m$ Linien aus, in denen $T=0$ ist, oder wenn man sie paarweise so
verbindet, dass jede, wo, bei wachsendem $\psi$, das Zeichen aus $-$ in $+$ übergeht, zusammen mit der nächstfolgenden, wo der entgegengesetzte Uebergang
Statt findet, wie die Begrenzungslinie eines Flächentheils mit positivem $T$ betrachtet wird, so treffen in $P$ überhaupt $m$ dergleichen Begrenzungslinien zusammen.

Von der andern Seite ist klar, dass so wie $Z$ unbestimmt durch $z^m$ und
durch keine höhere Potenz von $z$ theilbar ist, $X$ den Factor $(x-p)^m$, aber keine
höhere Potenz von $x-p$ enthalten wird. Es ist also allemal, wenn $p$ irgend
eine Wurzel der Gleichung $X=0$ bedeutet, der Exponent der höchsten Potenz von $x-p$, durch welche $X$ theilbar ist, der Anzahl der \emph{in} $P$ zusammentreffenden Begrenzungslinien für Flächen mit positivem $T$ gleich, oder was dasselbe ist, der Anzahl solcher \emph{an} $P$ zusammentreffender Flächen.

Uebrigens ist es leicht, der Beweisführung eine von Einmischung unendlich
kleiner Grössen ganz unabhängige Einkleidung zu geben, und zwar ganz analog
der Schlussreihe in den Art.\,3 und 4. Es lässt sich nemlich ein Werth von $s$ nachweisen, für welchen, so wie für jeden kleinern, der ganze Cyklus aller Werthe
von $\psi$ dieselbe abwechselnde Folge von $m$ Stücken mit positivem $T$ und ebensovielen mit negativem darbietet. Diese Eigenschaft hat die positive Wurzel der
Gleichung
\[
0=m\surd\frac{1}{2}-(m+1)ls-(m+2)l'ss-(m+3)l''s^3-\mbox{u.\,s.\,w.}
\]
wo $l$, $l'$, $l''$ u.\,s.\,w. die positiven Quadratwurzeln aus den Normen der complexen
Grössen $L$, $L'$, $L''$ u.\,s.\,w. bedeuten, oder wo
\[
\begin{array}{rl}
L&=l\,(\cos\lambda+i\sin\lambda)\\
L'&=l'(\cos\lambda'+i\sin\lambda')\\
L''&=l''(\cos\lambda''+i\sin\lambda'')\ \mbox{u.\,s.\,w.}
\end{array}
\]
gesetzt ist. Ich glaube jedoch, die sehr leichte Entwicklung dieses Satzes hier
übergehen zu können.

Schliesslich mag noch bemerkt werden, dass bei der Beweisführung in der
Abhandlung von 1799 die Betrachtung \emph{zweier} Systeme von Linien erforderlich
war, das eine die Linien wo $T=0$, das andere diejenigen wo $U=0$ enthaltend, während in unserm jetzigen Verfahren die Betrachtung Eines Systems
